%=====================================================================
\chapter{Reinforcement Learning}\label{ch:rl}
%=====================================================================
\locator{5}{Part V --- Sequences, Decisions and Reinforcement}

\begin{outcomes}
By the end of this chapter you will be able to:
\begin{tight}
  \item \textbf{Explain} how reinforcement learning differs from planning with a
        known model and from supervised learning.
  \item \textbf{Estimate} values by Monte Carlo averaging of observed returns.
  \item \textbf{Explain} the exploration--exploitation trade-off, and \textbf{apply}
        greedy, $\varepsilon$-greedy and UCB action selection to a bandit problem.
  \item \textbf{Perform} temporal-difference and Q-learning updates by hand.
  \item \textbf{Contrast} Q-learning with SARSA, and \textbf{describe} how function
        approximation leads to deep reinforcement learning.
\end{tight}
\end{outcomes}

\begin{prereq}
\chref{mdp} (MDPs, returns, value functions, the Bellman equations) and \chref{hmm}
(Monte Carlo estimation). \chref{generalization} for the trade-off between trying new
things and using what is known, which returns here in a new form.
\end{prereq}

\begin{keyterms}
reinforcement learning $\bullet$ model-free $\bullet$ episode $\bullet$ Monte Carlo
prediction $\bullet$ first-visit Monte Carlo $\bullet$ exploration vs.\ exploitation
$\bullet$ multi-armed bandit $\bullet$ greedy $\bullet$ $\varepsilon$-greedy
$\bullet$ upper confidence bound (UCB) $\bullet$ incremental mean $\bullet$ temporal
difference (TD) $\bullet$ TD error $\bullet$ bootstrapping $\bullet$ Q-learning
$\bullet$ SARSA $\bullet$ on-policy vs.\ off-policy $\bullet$ function
approximation $\bullet$ deep Q-network
\end{keyterms}

\section{Learning From Consequences}

\chref{mdp} solved decision problems when the rules of the world were written down:
every transition probability and every reward. \term{Reinforcement learning} (RL)
solves them when they are not. The agent acts, observes what happens and what it is
paid, and improves its behaviour from that experience alone. Nobody tells it the right
action --- unlike supervised learning, there are no labels, only a reward that may
arrive many steps after the action that earned it. That is how Samuel's checkers
program improved in 1959 (\chref{introduction}), and how AlphaGo beat a world champion
in 2016.

RL is \term{model-free} in the sense used here: it learns value functions or policies
directly, without first estimating $P(s' \mid s, a)$. Two ideas carry the whole
chapter. \textbf{Averaging}: the value of a state is the average return that follows
it, so average the returns you observe (Monte Carlo). \textbf{Bootstrapping}: the
Bellman equation says a state's value is the reward plus the next state's value, so
nudge your estimate towards that (temporal-difference learning).

\section{Monte Carlo Methods}

\chref{hmm} introduced \emph{Monte Carlo estimation}: estimate an expectation by
sampling and averaging, with error shrinking like $1/\sqrt{n}$. The simplest
illustration needs no MDP at all. In the model of \chref{hmm}, a healthy project
stays healthy for another week with probability 0.8; what is the chance it stays
healthy for ten weeks? Exactly, $0.8^{10} = 0.107$. By simulation: run the ten weeks
100 times and count the runs that stay healthy --- the lab gets 0.120; a million runs
give 0.1071. The same procedure works when no exact formula
exists, which is the point.

\begin{definitionbox}[Monte Carlo Prediction]
To estimate $V^{\pi}(s)$ without a model: run many \term{episodes} (complete
sequences from start to end) following $\pi$; each time $s$ is visited, record the
return $G$ that followed; estimate $V^{\pi}(s)$ as the average of the recorded returns.
\term{First-visit} Monte Carlo counts only the first visit to $s$ in each episode.
With an $\varepsilon$-greedy policy (below) and averaging $Q(s, a)$ instead of $V(s)$,
the same idea learns to \emph{control}, not just predict.
\end{definitionbox}

\begin{worked}[Averaging Returns]
Three episodes pass through state $s$, and the returns that followed the first visit
were $G = 7$, $3$ and $5$. The Monte Carlo estimate is $V(s) = (7 + 3 + 5)/3 =
\mathbf{5}$.

\smallskip
The average can be kept \emph{incrementally} --- no need to store every return:
\[
V_{n} = V_{n-1} + \frac{1}{n}\big(G_n - V_{n-1}\big)
\]
After the first episode $V = 7$; after the second $7 + \tfrac{1}{2}(3 - 7) = 5$; after
the third $5 + \tfrac{1}{3}(5 - 5) = 5$. The form ``old estimate plus a step towards
the new evidence'' appears in every method in this chapter.
\end{worked}

Monte Carlo is simple and unbiased, but it must wait until an episode ends to learn
anything, and returns summed over long episodes are noisy. The temporal-difference
methods later in the chapter fix both.

\section{Exploration and Exploitation}

An agent that always takes the action it currently believes best \emph{exploits}
what it knows, but may never discover that another action is better. One that tries
actions at random \emph{explores}, but wastes effort on actions it already knows are
poor. Every RL agent must balance the two, and the cleanest setting in which to study
the balance is the bandit.

\begin{definitionbox}[Multi-Armed Bandit]
$k$ actions (``arms'', after slot machines), each paying a random reward from an
unknown distribution; there is only one state. The agent pulls one arm per step and
tries to maximise its total reward. Its estimate $Q(a)$ of each arm's value is the
average reward received from it so far, updated incrementally:
$Q(a) \leftarrow Q(a) + \tfrac{1}{N(a)}\big(r - Q(a)\big)$.
\end{definitionbox}

\rowcolors{2}{white}{lightbg}
\begin{longtable}{@{}L{2.8cm}L{6.4cm}L{5.8cm}@{}}
\toprule
\rowcolor{primary!15}
\textbf{Strategy} & \textbf{Rule} & \textbf{Behaviour} \\
\midrule
\endhead
Greedy & Always pull $\arg\max_a Q(a)$ & Locks onto whichever arm looked good first;
may never find the best \\
$\varepsilon$-greedy & With probability $\varepsilon$ pull a random arm; otherwise
greedy & Keeps exploring forever at rate $\varepsilon$; simple and robust \\
UCB & Pull $\arg\max_a \Big[Q(a) + c\sqrt{\dfrac{\ln t}{N(a)}}\Big]$ & Explores arms in
proportion to how \emph{uncertain} their estimate is: optimism under uncertainty \\
\bottomrule
\end{longtable}

\begin{worked}[Which Estimation Technique Next Sprint?]
A team is trying three estimation techniques --- A planning poker, B t-shirt sizing,
C three-point estimates --- one per sprint. The reward is 1 when the sprint's
estimate lands within 10\% of the actual effort, 0 otherwise. So far the rewards
average: A $0.5$ (used in 10 sprints), B $0.6$ (5 sprints), C $0.3$ (2 sprints).
Total pulls $t = 17$; $\ln 17 = 2.833$.

\smallskip
\textbf{Greedy:} B, the highest average.

\textbf{$\varepsilon$-greedy} ($\varepsilon = 0.1$): B with probability
$0.9 + 0.1/3 = 0.933$; A and C with $0.033$ each.

\textbf{UCB} ($c = 1$):
\begin{tight}
  \item A: $0.5 + \sqrt{2.833/10} = 0.5 + 0.532 = 1.032$
  \item B: $0.6 + \sqrt{2.833/5} = 0.6 + 0.753 = 1.353$
  \item C: $0.3 + \sqrt{2.833/2} = 0.3 + 1.190 = \mathbf{1.490}$
\end{tight}
UCB chooses \textbf{C}: its average is lowest, but it has been tried only twice, so
its true value is the most uncertain. If C really is poor, a few more sprints will
lower its bonus and UCB will stop choosing it; if it is good, greedy would never have
found out.

\smallskip
\textbf{Update} after B is used again and its estimate lands (reward 1): $N(B) = 6$
and $Q(B) = 0.6 + \tfrac{1}{6}(1 - 0.6) = 0.667$.
\end{worked}

\dsfig{%
\begin{tikzpicture}
\begin{axis}[width=10cm, height=5.2cm, axis lines=left,
  xlabel={step}, ylabel={average reward (20-step window)},
  tick label style={font=\tiny}, label style={font=\scriptsize},
  xmin=0, xmax=1000, ymin=0.3, ymax=1.65,
  legend style={font=\tiny, draw=gray!40, at={(0.98,0.03)}, anchor=south east}]
\addplot[gray!60!black, line width=1.2pt] coordinates {(20,0.743) (29,0.956) (69,1.042) (109,1.060) (149,1.072) (189,1.055) (229,1.068) (269,1.083) (309,1.078) (349,1.067) (389,1.059) (429,1.072) (469,1.060) (509,1.066) (549,1.064) (589,1.060) (629,1.070) (669,1.059) (709,1.064) (749,1.059) (789,1.063) (829,1.056) (869,1.067) (909,1.058) (949,1.076) (989,1.077) (1000,1.079)};
\addlegendentry{greedy}
\addplot[secondary, line width=1.2pt] coordinates {(20,0.738) (29,0.960) (69,1.091) (109,1.184) (149,1.234) (189,1.281) (229,1.301) (269,1.326) (309,1.319) (349,1.322) (389,1.349) (429,1.361) (469,1.354) (509,1.355) (549,1.374) (589,1.391) (629,1.369) (669,1.369) (709,1.385) (749,1.380) (789,1.364) (829,1.380) (869,1.368) (909,1.376) (949,1.378) (989,1.388) (1000,1.385)};
\addlegendentry{$\varepsilon$-greedy, $\varepsilon = 0.1$}
\addplot[accent, line width=1.2pt] coordinates {(20,0.417) (29,0.799) (69,1.082) (109,1.237) (149,1.325) (189,1.325) (229,1.348) (269,1.384) (309,1.419) (349,1.399) (389,1.425) (429,1.418) (469,1.439) (509,1.444) (549,1.445) (589,1.458) (629,1.472) (669,1.467) (709,1.453) (749,1.476) (789,1.462) (829,1.453) (869,1.471) (909,1.469) (949,1.496) (989,1.479) (1000,1.464)};
\addlegendentry{UCB, $c = 2$}
\addplot[greenhl!70!black, dashed] coordinates {(0,1.54)(1000,1.54)};
\node[font=\tiny, greenhl!45!black, anchor=south west] at (axis cs:10,1.55) {best arm, on average};
\end{axis}
\end{tikzpicture}}%
{Ten-armed bandits (600 random problems, averaged). Greedy settles early on a
mediocre arm; $\varepsilon$-greedy keeps improving; UCB explores hardest at first --- its
worst start --- and earns the most thereafter.}{bandits}

\section{Temporal-Difference Learning}

Monte Carlo waits for the end of an episode to see the return. \term{Temporal-difference}
(TD) learning updates after every single step, using the Bellman equation: the value
of a state should equal the reward plus the discounted value of the next state, so move
the estimate towards that.

\begin{definitionbox}[TD(0) Update]
After moving from $s$ to $s'$ with reward $r$:
\[
V(s) \leftarrow V(s) + \alpha\big[\underbrace{r + \gamma V(s')}_{\text{TD target}} - V(s)\big]
\]
The bracket is the \term{TD error}: how much better or worse things turned out than
expected. Using the current estimate $V(s')$ inside the target is called
\term{bootstrapping} --- learning a guess from a guess. It lets TD learn during an
episode, and in tasks that never end.
\end{definitionbox}

\begin{worked}[One TD Update]
$V(s) = 5$ and $V(s') = 10$. The agent moves from $s$ to $s'$ and receives $r = 2$;
$\gamma = 0.9$, $\alpha = 0.1$.
\[
\text{TD target} = 2 + 0.9 \times 10 = 11, \quad \text{TD error} = 11 - 5 = 6, \quad
V(s) \leftarrow 5 + 0.1 \times 6 = \mathbf{5.6}
\]
The step was better than expected, so the estimate of $s$ rises --- by a tenth of the
surprise, because one step is weak evidence.
\end{worked}

\section{Q-Learning}

To choose actions without a model, the agent needs action values $Q(s, a)$, not state
values: with $Q$, the best action is simply $\arg\max_a Q(s, a)$, with no need to know
where each action leads.

\begin{definitionbox}[Q-Learning (Watkins, 1989)]
Initialise $Q(s, a)$ arbitrarily (say 0). In each step: in state $s$ choose $a$
(e.g.\ $\varepsilon$-greedily from $Q$), observe reward $r$ and next state $s'$, and
update
\[
Q(s, a) \leftarrow Q(s, a) + \alpha\Big[r + \gamma\max_{a'} Q(s', a') - Q(s, a)\Big]
\]
Under mild conditions --- every action tried in every state infinitely often, and a
suitably decreasing $\alpha$ --- $Q$ converges to the optimal $Q^{*}$ of \chref{mdp},
whatever exploratory policy generated the experience.
\end{definitionbox}

\begin{worked}[Three Q-Learning Updates on the Technical-Debt MDP]
The agent does not know the transition probabilities of \dsref{debt-mdp}; it only
experiences them, one sprint at a time. $Q$ starts at 0 everywhere; $\alpha = 0.5$,
$\gamma = 0.9$.

\smallskip
\textbf{Step 1:} in H, ship; reward $+10$; stays in H.
$Q(H, \text{ship}) \leftarrow 0 + 0.5\,[10 + 0.9 \times \max(0, 0) - 0] = \mathbf{5}$.

\textbf{Step 2:} in H, ship; reward $+10$; debt builds up, now D.
$Q(H, \text{ship}) \leftarrow 5 + 0.5\,[10 + 0.9 \times \max(Q(D, \cdot)) - 5] = 5 + 0.5 \times 5 =
\mathbf{7.5}$.

\textbf{Step 3:} in D, refactor; reward $-4$; back to H.
$Q(D, \text{ref}) \leftarrow 0 + 0.5\,[-4 + 0.9 \times 7.5 - 0] = 0.5 \times 2.75 =
\mathbf{1.375}$.

\smallskip
The refactoring was penalised $-4$, yet its Q-value is \emph{positive}: the update
already credits it with the value of the healthy state it leads to. Run for 200{,}000
steps (the lab), Q-learning arrives at the same policy value iteration computed from
the full model --- ship when healthy, refactor when debt-laden or frozen --- having
never seen a single transition probability.
\end{worked}

\begin{conceptbox}[Q-Learning and SARSA: Off-Policy and On-Policy]
Q-learning's target uses $\max_{a'} Q(s', a')$ --- the value of the \emph{best} next
action, whether or not the agent will take it. It learns the value of the optimal
policy while behaving exploratorily: \term{off-policy}. \textbf{SARSA} uses instead the
action $a'$ the agent actually takes next,
$Q(s, a) \leftarrow Q(s, a) + \alpha[r + \gamma Q(s', a') - Q(s, a)]$ (the name lists the
quintuple $s, a, r, s', a'$). It learns the value of the policy it follows, exploration
included: \term{on-policy}. Where an exploratory step is disastrous --- pushing an
untested build to production, say --- SARSA learns a safer policy than Q-learning,
because it accounts for its own occasional random steps.
\end{conceptbox}

\section{From Tables to Deep Reinforcement Learning}

A table with one entry per state--action pair is fine for three codebase states and
hopeless for the full state of a project portfolio --- every team, task and deadline.
\term{Function approximation} replaces the table by a model $\hat{Q}(s, a; \mathbf{w})$
--- a linear model over state features, or a neural network (\chref{ann}) ---
trained with the same TD targets as regression targets. A deep neural network
trained this way is a \term{deep Q-network} (DQN), and the field is \term{deep
reinforcement learning}, the technology behind AlphaGo. Approximation
brings instability that the tabular case does not have, which is why deep RL systems
need engineering tricks (experience replay, target networks) that belong to a deep
RL course.

\begin{alertbox}[What Makes Real Reinforcement Learning Hard]
\begin{tight}
  \item \textbf{Samples are expensive.} Q-learning needed 200{,}000 steps on a
        three-state problem; a real team runs perhaps twenty-five sprints a year.
        Most practical RL learns in a simulator first.
  \item \textbf{Exploration can be unsafe.} ``Try a random action'' is unacceptable
        for a production release or a client's budget.
  \item \textbf{The reward is the specification.} An agent will exploit any gap
        between the reward and what was meant.
\end{tight}
\end{alertbox}

\begin{pitfall}
\begin{tight}
  \item \textbf{Never exploring.} A greedy agent can lock onto a poor action for
        ever.
  \item \textbf{Exploring for ever at a high rate.} Decay $\varepsilon$ over time once
        estimates are reliable.
  \item \textbf{Confusing Q-learning's target with SARSA's.} Max over next actions
        versus the next action actually taken.
  \item \textbf{A learning rate that is too large.} Values oscillate; too small and
        learning takes for ever.
  \item \textbf{Judging a policy by one run.} Rewards are random; average many
        episodes, as in the bandit figure.
\end{tight}
\end{pitfall}

%---------------------------------------------------------------------
\begin{lab}[Monte Carlo, Bandits and Q-Learning]
\begin{lstlisting}[language=Python]
import numpy as np

rng = np.random.default_rng(0)

# --- 1. Monte Carlo inference: estimate an expectation by averaging ------
# the chance that a project stays healthy for 10 weeks (it slips with
# probability 0.2 each week): exact answer 0.8 ** 10 = 0.107
for n in (100, 10_000, 1_000_000):
    runs = rng.random((n, 10)) > 0.2            # True = stayed healthy
    print(f"{n:9,d} simulated runs: P(healthy 10 weeks) ~ "
          f"{runs.all(axis=1).mean():.4f}")


# --- 2. exploration vs exploitation: a 10-armed bandit -------------------
def bandit(strategy, steps=1000, arms=10, runs=300):
    total = np.zeros(steps)
    for _ in range(runs):
        true = rng.normal(0, 1, arms)           # unknown mean payoffs
        Q, N = np.zeros(arms), np.zeros(arms)
        for t in range(1, steps + 1):
            if strategy == "greedy":
                a = int(np.argmax(Q))
            elif strategy == "e-greedy":
                a = (int(rng.integers(arms)) if rng.random() < 0.1
                     else int(np.argmax(Q)))
            else:                               # UCB, c = 2
                ucb = Q + 2 * np.sqrt(np.log(t) / np.maximum(N, 1e-9))
                a = int(np.argmax(np.where(N == 0, np.inf, ucb)))
            r = rng.normal(true[a], 1)
            N[a] += 1
            Q[a] += (r - Q[a]) / N[a]           # incremental average
            total[t - 1] += r
    return total / runs


for s in ("greedy", "e-greedy", "ucb"):
    avg = bandit(s)
    print(f"{s:8s}: mean reward, steps 1-100 {avg[:100].mean():.2f}; "
          f"steps 901-1000 {avg[900:].mean():.2f}")

# --- 3. Q-learning on the technical-debt MDP, without being told P -------
P = {("H", "ship"): ([0.8, 0.2, 0.0], 10),
     ("D", "ship"): ([0.0, 0.5, 0.5], 5),
     ("F", "ship"): ([0.0, 0.0, 1.0], 0),
     ("H", "refactor"): ([1, 0, 0], -2), ("D", "refactor"): ([1, 0, 0], -4),
     ("F", "refactor"): ([1, 0, 0], -15)}
S, A = ["H", "D", "F"], ["ship", "refactor"]


def env_step(s, a):                         # the agent only sees samples
    probs, reward = P[(S[s], A[a])]
    return int(rng.choice(3, p=probs)), reward


Q = np.zeros((3, 2))
alpha, gamma, eps = 0.1, 0.9, 0.1
s = 0
for t in range(200_000):
    a = int(rng.integers(2)) if rng.random() < eps else int(Q[s].argmax())
    s2, r = env_step(s, a)
    Q[s, a] += alpha * (r + gamma * Q[s2].max() - Q[s, a])   # Q-learning
    s = s2
print("learned Q (rows H, D, F; columns ship, refactor):")
print(Q.round(1))
print("learned policy:", dict(zip(S, [A[a] for a in Q.argmax(axis=1)])))
\end{lstlisting}
The learned Q-values approach the true $Q^{*}$ (for H: ship 78.6, refactor 68.8) but
do not reach them exactly, because $\alpha$ stays at 0.1 and the agent keeps
exploring. Try $\varepsilon = 0$: does Q-learning still discover that refactoring
pays?
\end{lab}

%---------------------------------------------------------------------
\begin{checkpoint}
\begin{tightnum}
  \item Returns of 10, 4 and 7 followed three visits to a state. What is the Monte
        Carlo estimate of its value, and what does it become after a fourth return of
        3?
  \item Why can a purely greedy bandit strategy perform badly for ever?
  \item $Q(s, a) = 4$; the agent receives $r = 1$ and lands in $s'$ where the best
        action value is 6. With $\alpha = 0.5$ and $\gamma = 0.9$, what is the new
        $Q(s, a)$?
\end{tightnum}
\end{checkpoint}

%---------------------------------------------------------------------
\begin{chaptersummary}
\begin{tight}
  \item Reinforcement learning finds good policies from experience, without a model
        of the environment and without labels.
  \item Monte Carlo methods average the returns observed after visiting a state;
        they are unbiased but wait for episodes to end.
  \item Exploration versus exploitation: greedy action selection can lock in a poor
        choice; $\varepsilon$-greedy explores at a fixed rate; UCB explores where
        estimates are most uncertain.
  \item TD learning updates after every step towards $r + \gamma V(s')$,
        bootstrapping from its own estimates.
  \item Q-learning updates action values towards $r + \gamma\max_{a'}Q(s', a')$ and
        converges to $Q^{*}$ off-policy; SARSA is its on-policy counterpart.
  \item Function approximation, with neural networks, extends RL to huge state
        spaces: deep reinforcement learning.
\end{tight}
\end{chaptersummary}

\begin{reviewq}
  \item \textbf{(MCQ)} Q-learning is: (a)~on-policy (b)~off-policy (c)~a supervised
        method (d)~a planning method that needs the model.
  \item \textbf{(MCQ)} In UCB, an arm that has been tried few times receives:
        (a)~a smaller bonus (b)~a larger bonus (c)~no bonus (d)~a negative value.
  \item \textbf{(MCQ)} The TD error is: (a)~$r - V(s)$ (b)~$r + \gamma V(s') - V(s)$
        (c)~$V(s') - V(s)$ (d)~$G - r$.
  \item \textbf{(Short)} Compare Monte Carlo and TD learning: when does each update,
        and what does each target use?
  \item \textbf{(Short)} Explain ``optimism under uncertainty'' with reference to UCB.
  \item \textbf{(Short)} Why is exploration dangerous in some real applications, and
        what is done about it?
  \item \textbf{(Applied)} For the estimation example, compute the UCB scores with
        $c = 2$ and state which technique is chosen. Then suppose it is used and the
        estimate misses (reward 0): recompute its average.
  \item \textbf{(Applied)} Continue the Q-learning worked example with step 4: in H,
        refactor, reward $-2$, stays in H. Compute the new $Q(H, \text{refactor})$
        and state which action is now greedy in H.
\end{reviewq}
